\documentclass[12pt, a4paper]{article}
\usepackage{etex} % расширение классического tex в частности позволяет подгружать гораздо больше пакетов, чем мы и займёмся далее

%%%%%%%%%% Математика %%%%%%%%%%
\usepackage{amsmath,amsfonts,amssymb,amsthm,mathtools} 
%\mathtoolsset{showonlyrefs=true}  % Показывать номера только у тех формул, на которые есть \eqref{} в тексте.
%\usepackage{leqno} % Нумерация формул слева


%%%%%%%%%%%%%%%%%%%%%%%% Шрифты %%%%%%%%%%%%%%%%%%%%%%%%%%%%%%%%%
\usepackage{fontspec}         % пакет для подгрузки шрифтов
\setmainfont{Arial}   % задаёт основной шрифт документа

\defaultfontfeatures{Mapping=tex-text}

\newfontfamily{\cyrillicfonttt}{Arial}
\newfontfamily{\cyrillicfont}{Arial}
\newfontfamily{\cyrillicfontsf}{Arial}

\usepackage{unicode-math}     % пакет для установки математического шрифта
\setmathfont{Asana Math}      % шрифт для математики

\usepackage{polyglossia}      % Пакет, который позволяет подгружать русские буквы
\setdefaultlanguage{russian}  % Основной язык документа
\setotherlanguage{english}    % Второстепенный язык документа


%%%%%%%%%% Работа с картинками %%%%%%%%%
\usepackage{graphicx}                  % Для вставки рисунков
\usepackage{graphics}
\graphicspath{{images/}{pictures/}}    % можно указать папки с картинками
\usepackage{wrapfig}                   % Обтекание рисунков и таблиц текстом


%%%%%%%%%% Графика и рисование %%%%%%%%%%
\usepackage{tikz, pgfplots}  % язык для рисования графики из latex'a
\usepackage{amscd}                  %Пакеты для рисования
\usepackage[matrix,arrow,curve]{xy} %комунитативных диаграмм


%%%%%%%%%% Теоремы %%%%%%%%%%
\theoremstyle{plain}              % Это стиль по умолчанию.  Есть другие стили. 
\newtheorem{theorem}{Теорема}[section]
\newtheorem{result}{Следствие}[theorem]
% счётчик подчиняется теоремному, нумерация идёт по главам согласованно между собой

\theoremstyle{definition}         % убирает курсив и что-то еще наверное делает ;)
\newtheorem*{defin}{Определение}  % нумерация не идёт вообще

\newtheorem{fignia}{Какая-то фигня}


%%%%%%%%%% Свои команды %%%%%%%%%%
\usepackage{etoolbox}    % логические операторы для своих макросов

% Пакет, который ставит в каждом первом абзаце главы красную строку
% Просто, чтобы эта pdf-ка нормально смотрелась :)
\usepackage{indentfirst}
\setkeys{russian}{babelshorthands=true}


\DeclareMathOperator{\Var}{Var}
\DeclareMathOperator{\Cov}{Cov}


\begin{document}

\section{Часть 1}

\[\Var(X) \qquad \text{VS} \qquad Var(X) \]
\[\Cov(X,Y) \qquad \text{VS} \qquad Cov(X,Y) \]

\begin {center}

\def\s{\ensuremath{\sigma}}
$\sigma$-алгебра \quad VS \quad \s-алгебра 


\newcommand{\co}{x\_1  \backslash ldots x\_n}
\co

\newcommand{\com}[2]{x\_#1  \backslash ldots x\_#2}
\com{a}{z} \\
\com{1}{6} \\
\com{(a,b)}{(c,d)}

\end{center}


\def\tk{
\begin{tikzpicture}
\draw[blue,fill=blue](2,2) circle [radius = 0.1];
\end{tikzpicture}
}

\begin {itemize}
\item [\tk] Первый пункт
\item [\tk] Второй пункт
\item [\tk] Третий пункт
\end {itemize}


\newcommand{\llim}[2]{\lim\limits_{n \to #1} #2}
$\llim {\infty} {\frac{n+1}{n}}$


\renewcommand\thefigure{\thesection.\arabic{figure}}    
\setcounter{figure}{0}   
 
\begin{figure}[h!]
\begin{center}
\includegraphics[width=0.3\textheight]{doge.png}
\caption{Хаски}
\end{center}
\end{figure}


\section{Часть 2. Картиночки}
\setcounter{figure}{0}    
\begin{figure}[h!]
\begin{center}
\includegraphics[width=0.4\textheight]{doge.png}
\caption{Больше хаски}
\end{center}
\end{figure}
  
\begin{figure}[h!]
\begin{center}
\includegraphics[width=0.5\textheight]{doge.png}
\caption{Еще больше хаски}
\end{center}
\end{figure}

\newcounter{i}
\newcommand{\eq}{%
\addtocounter{i}{1}   
Eq. (\arabic{i})\\}

\begin{equation}
    E=mc^2\tag{\eq}
  \end{equation}
    
\end{document}